\begin{qfmethod}[ESG considerations in managing long-term institutional portfolios: motives, approaches from exclusion to stewardship, costs of constraints, climate risk and measurement]\label{meth:institutional-esg-integration}
\qffield{Statement} \emph{Why long-term institutions consider ESG}: as universal owners of diversified portfolios they bear system-wide risks such as climate change; their mandates may require it (Norway's mandate expects portfolio companies to align with net zero by 2050, \citep[climate action plan]{nbim-gpfg-2026}); and beneficiaries' preferences confer legitimacy \citep[chap.~3, pp.~71--73, 112]{s7-ang}.

\emph{Approaches}, as defined by CFA Institute, GSIA and PRI: screening (negative, norms-based, best-in-class), ESG integration into analysis, thematic investing, stewardship (engagement and voting) and impact investing \citep[Tables~1--2, p.~7]{sustainable-investment-review-2022}. Norway combines engagement with exclusions recommended by an independent Council on Ethics, as with Wal-Mart in 2006 (Ang, pp.~71--73).

\emph{Costs and evidence}: exclusions shrink diversification, but excluding a few sectors from a diversified portfolio costs little in mean--variance terms; screening on characteristics is an active strategy judged like any other, and the evidence on returns is mixed (pp.~110--112). Fiduciary duty requires that non-financial aims be disclosed (GAPP~19.1 \citep[GAPP~19.1]{santiago-principles-2008}) and that the cost be measured against the full opportunity set (Ang, p.~112).

\emph{Measurement}: tracking error from exclusions, weighted average carbon intensity, scenario analysis of transition and physical risk, and engagement outcomes.

\qffield{Practice} (original; \nolinkurl{trading_notes.numerical.institutional_portfolios.institutional_portfolios_checks}, part~5).
\begin{itemize}
\item \emph{Exclusion cost}: removing a sector with 1.5 percent benchmark weight whose return differs from the rest with 25 percent volatility adds about 0.375 percent of tracking error; four sectors totalling 5.5 percent add 0.72 percent, within a 1.25 percent budget.
\item \emph{Carbon}: holdings of 40/35/25 percent with intensities of 20, 150 and 600 tonnes per USD~million of revenue give a weighted intensity of 210.5; moving 10 points from the most to the least intensive cuts it to 152.5.
\item \emph{Recommendation for the Elmsworth endowment}: adopt a policy based on engagement and voting through its managers, a short norms-based exclusion list approved by the board, climate scenarios in the annual stress test, and reporting of tracking error from exclusions; avoid broad exclusions that would push tracking error beyond its budget without a stated mandate.
\end{itemize}

\qffield{Conditions} The carbon and tracking-error figures are illustrative; ESG ratings differ across providers.

\qffield{Pitfalls} Treating ESG integration as a guarantee of higher returns. Excluding without measuring the cost.

\qffield{See also} \hyperref[meth:esg-exposure-evaluation]{ESG exposures}; \hyperref[meth:swf-risk-case]{SWF case}.

\end{qfmethod}
